\documentclass[11pt,a4paper]{article}
\usepackage[margin=1in]{geometry}
\usepackage[T1]{fontenc}
\usepackage{lmodern}
\usepackage{booktabs}
\usepackage{array}
\usepackage{longtable}
\usepackage{enumitem}
\usepackage[hidelinks]{hyperref}
\setlist[itemize]{leftmargin=1.4em}
\setlist[enumerate]{leftmargin=1.6em}

\title{\textbf{5G Innovation Hackathon 2026 Proposal}\\
\large Gesture-to-Voice over 5G: Real-Time Assistive Communication Platform}
\author{Team Name: [Enter Team Name] \\
Institute/Organization: [Enter Institute/Startup/Organization] \\
Primary Contact: [Name, Email, Phone]}
\date{Submission Date: [DD Month 2026]}

\begin{document}
\maketitle

\section*{1. Executive Summary}
This proposal presents an assistive communication solution that converts hand gestures into spoken output in real time and extends it into a 5G-native service for low-latency, scalable deployment. Our base prototype (already implemented) performs webcam-based hand landmark extraction, gesture classification, debounced sentence construction, and asynchronous text-to-speech output. 

In the hackathon, we will evolve this into a robust 5G use case by integrating edge inference, QoS-aware streaming, and optional cloud services for multilingual speech and remote caregiver dashboards. The system targets inclusive communication in hospitals, public-service counters, classrooms, and disaster-response settings where rapid non-verbal-to-verbal translation can improve accessibility and safety.

\section*{2. Problem Statement}
People with speech and hearing disabilities frequently face communication barriers in time-sensitive situations, especially when interpreters are unavailable. Existing mobile accessibility tools often struggle with:
\begin{itemize}
    \item latency and jitter in real-time translation,
    \item poor reliability in variable network conditions,
    \item limited vocabulary and low contextual personalization,
    \item low interoperability with healthcare, emergency, and enterprise systems.
\end{itemize}

A 5G-enabled architecture can address these issues through low latency, improved reliability, network slicing, and edge computing support.

\section*{3. Proposed Solution}
\textbf{Project Name:} GestureVoice-5G

\textbf{Core Idea:} Capture hand gestures on camera-enabled devices, infer intended tokens (words/characters), build context-aware text, and generate real-time speech output, while using 5G network capabilities for performance guarantees.

\subsection*{3.1 Current Working Prototype (Completed)}
Our existing implementation already includes:
\begin{itemize}
    \item Real-time hand tracking using MediaPipe Hands (21 landmarks).
    \item Wrist-centered, scale-invariant landmark normalization.
    \item Dual classifier path: trained MLP model and rule-based fallback.
    \item Stable multi-frame debouncing and cooldown logic for robust token commit.
    \item Word mode and character mode with control gestures (space, delete, clear, speak).
    \item Non-blocking asynchronous TTS queue.
    \item Multiple TTS backends with fallback (Kokoro/Coqui/pyttsx3/macOS \texttt{say}).
    \item UI overlays for confidence, FPS, active mode, and output sentence.
\end{itemize}

\subsection*{3.2 5G Hackathon Enhancements (Planned)}
\begin{itemize}
    \item Edge deployment of inference service in 5G lab MEC node.
    \item Adaptive offloading: on-device first, edge fallback for heavy models.
    \item QoS-aware streaming profile for low-latency gesture frames and speech packets.
    \item Network slicing policy experiments for assistive traffic prioritization.
    \item IoT integration APIs for nurse-station displays, classroom boards, and emergency consoles.
    \item Multilingual TTS/translation pipeline for Indian language support.
    \item Secure event logging and analytics dashboard for service quality monitoring.
\end{itemize}

\section*{4. Alignment with Hackathon Thematic Areas}
\begin{longtable}{p{0.35\linewidth} p{0.6\linewidth}}
\toprule
\textbf{Thematic Area} & \textbf{Alignment in This Proposal} \\
\midrule
5G/5G Advanced use cases & Real-time assistive communication over 5G with edge inference and reliability testing. \\
AI/ML for telecom technology & ML-based gesture recognition with confidence filtering and adaptive inference placement. \\
IoT-enabled applications (healthcare/education/industry) & Integration with healthcare counters, classrooms, and public service kiosks via API/IoT endpoints. \\
Disaster management / RedCap / FWA & Fast non-verbal communication support in emergency shelters and field operations; lightweight client profile for constrained devices. \\
NTN integration potential & Future-ready architecture supports delayed/remote inference fallback through satellite-backed connectivity scenarios. \\
Advanced 5G features (slicing, QoS) & Dedicated assistive communication slice experiments and QoS profile benchmarking. \\
Telecom security & Secure transport, access control, anonymized telemetry, and compliance-oriented logging. \\
\bottomrule
\end{longtable}

\section*{5. System Architecture}
\textbf{Pipeline:}
\begin{enumerate}
    \item \textbf{Capture Layer:} Smartphone/webcam stream from user endpoint.
    \item \textbf{Perception Layer:} Hand landmark extraction and normalization.
    \item \textbf{Inference Layer:} Gesture classification (local/edge model selection).
    \item \textbf{Language Layer:} Debounced token aggregation and sentence generation.
    \item \textbf{Voice Layer:} Low-latency multilingual TTS output.
    \item \textbf{Service Layer:} Remote dashboard, alerting, analytics, and logs.
\end{enumerate}

\textbf{5G Role:}
\begin{itemize}
    \item Uplink for real-time gesture frames/feature vectors.
    \item Edge-hosted inference for reduced round-trip delay.
    \item Downlink for speech output and service updates.
    \item QoS/slicing controls for predictable assistive service quality.
\end{itemize}

\section*{6. Novelty and Differentiation}
\begin{itemize}
    \item Practical hybrid inference strategy: deterministic fallback + trainable model.
    \item Robust interaction design using multi-frame stability and cooldown control.
    \item 5G-native quality engineering instead of best-effort internet operation.
    \item Deployment-focused design for real service environments, not only lab demos.
\end{itemize}

\section*{7. Implementation Plan and Milestones}
\begin{longtable}{p{0.12\linewidth} p{0.32\linewidth} p{0.46\linewidth}}
\toprule
\textbf{Phase} & \textbf{Duration} & \textbf{Deliverables} \\
\midrule
P1 & Week 1--2 & Baseline hardening, dataset expansion, multilingual phrase set definition. \\
P2 & Week 3--4 & Edge inference containerization and 5G lab integration setup. \\
P3 & Week 5--6 & QoS/slicing experiments, latency benchmarking, packet-loss resilience tests. \\
P4 & Week 7--8 & Healthcare/classroom IoT integration, dashboard MVP, role-based access controls. \\
P5 & Week 9--10 & Pilot demonstrations in target scenarios; user feedback loop. \\
P6 & Week 11--12 & Final optimization, documentation, commercialization and IPR dossier. \\
\bottomrule
\end{longtable}

\section*{8. Expected Outcomes and Impact}
\subsection*{8.1 Technical Targets}
\begin{itemize}
    \item End-to-end response time target: sub-250 ms (gesture-to-audio in edge mode).
    \item Gesture commit stability: $\geq$ 95\% for trained vocabulary in controlled settings.
    \item Availability target: $\geq$ 99\% session continuity under tested 5G lab conditions.
\end{itemize}

\subsection*{8.2 Societal Impact}
\begin{itemize}
    \item Improves communication accessibility for speech/hearing-impaired users.
    \item Reduces dependency on human intermediaries in critical contexts.
    \item Enables inclusive public-service delivery and assistive digital infrastructure.
\end{itemize}

\section*{9. Scalability and Market Readiness}
\begin{itemize}
    \item \textbf{Primary sectors:} Hospitals, education institutes, public-service kiosks, enterprise support desks.
    \item \textbf{Deployment model:} SaaS + edge appliance + on-prem enterprise options.
    \item \textbf{Business model:} Subscription licensing for institutions, device bundle partnerships, API usage plans.
    \item \textbf{Go-to-market:} Pilot with institute 5G lab partners, then state-level accessibility and smart-city programs.
\end{itemize}

\section*{10. Risk Assessment and Mitigation}
\begin{longtable}{p{0.34\linewidth} p{0.6\linewidth}}
\toprule
\textbf{Risk} & \textbf{Mitigation Strategy} \\
\midrule
Gesture variability across users & User-specific calibration flow + mixed dataset training + fallback rules. \\
Network fluctuation & Adaptive mode switching (on-device/edge), local cache, QoS tuning. \\
Low-light or occlusion errors & Data augmentation, confidence gating, user guidance overlays. \\
Privacy/security concerns & Role-based access, encryption in transit, anonymized logs, audit trail. \\
Language diversity challenge & Modular language packs and phased multilingual rollout. \\
\bottomrule
\end{longtable}

\section*{11. IPR, Ethics, and Compliance}
\begin{itemize}
    \item File patentability assessment for adaptive gesture-to-voice over QoS-managed 5G architecture.
    \item Maintain ownership of in-house models, pipeline orchestration, and deployment workflows.
    \item Apply privacy-by-design principles and store only minimal telemetry needed for reliability improvement.
    \item Ensure inclusive dataset and UI testing across gender, age group, and skin-tone diversity where feasible.
\end{itemize}

\section*{12. Seed Funding Utilization Plan}
Requested seed funding will be used for:
\begin{itemize}
    \item data collection and annotation support,
    \item edge deployment and testing infrastructure,
    \item multilingual TTS integration and validation,
    \item pilot deployment logistics and user studies,
    \item IPR filing and compliance documentation.
\end{itemize}

\section*{13. Team and Collaboration}
\begin{itemize}
    \item Team Composition: Undergraduate/Postgraduate Students + Mentor + Optional Startup/Professional Collaborator.
    \item Internal Roles: AI/ML lead, telecom systems lead, frontend/integration lead, field testing lead.
    \item External Support: 5G lab mentors, industry advisors, accessibility-domain experts.
\end{itemize}

\section*{14. Conclusion}
GestureVoice-5G is an implementation-ready, socially relevant, and technically scalable 5G use case. We already have a functioning real-time gesture-to-voice baseline and a clear plan to convert it into a production-grade 5G assistive platform. The project aligns with the hackathon's focus on innovation, technical execution, market readiness, and societal impact.

\vspace{1em}
\noindent\textbf{Hackathon Timeline Note:} This proposal is prepared for the 5G Innovation Hackathon 2026 call window (20 March 2026 to 17 April 2026).

\end{document}
